\subsection{Retractions and contractions}

DEFINITION 4. --- Let X be a topological space and let A be a subset of X.

We call a retraction of X onto A a continuous map from X to A which is a retraction of the canonical injection of A into X. If there exists a retraction of X onto A, we say that X can be retracted onto A, or again that A is a retract of X.

Let X be a topological space and let A be a subspace of X.

Suppose there exists a retraction \(r\) of X onto A. The subspace A is identified with the set of points \(x \in \mathrm{X}\) such that \(x = r(x)\). If X is a separated space, then A is closed in X.

By the following lemma, for a subspace A to be a retract of X, it is necessary and sufficient that every continuous map from A to a topological space Y extends to a continuous map from X to Y.

Lemma 1. --- Let X and Y be topological spaces, and let \(f\colon X\to Y\) be a continuous map. The following conditions are equivalent:

\begin{enumerate}
\def\labelenumi{(\roman{enumi})}
\item
  For every topological space Z and every continuous map \(g: X \to Z\), there exists a continuous map \(g': Y \to Z\) such that \(g = g' \circ f\);
\item
  The map f is injective and has a continuous retraction;
\item
  The map f defines a homeomorphism from X onto its image \(f(\mathrm{X})\), which is a retract of Y.
\end{enumerate}

Suppose that \(f\) is injective, and let \(r: \mathrm{Y} \to \mathrm{X}\) be a continuous retraction of the map \(f\). For every continuous map \(g: \mathrm{X} \to \mathrm{Z}\), the map \(g' = g \circ r: \mathrm{Y} \to \mathrm{Z}\) satisfies \(g' \circ f = g \circ r \circ f = g\). This proves that (ii)⇒(i).

Suppose that condition (i) is satisfied. Let \(g \colon X \to X\) be the identity map. By hypothesis, there exists a continuous map \(g' \colon Y \to X\) such that \(g' \circ f = g\). This implies that the map \(f\) is injective and that \(g'\) is a continuous retraction of \(f\).

The equivalence of properties (ii) and (iii) is immediate.

DEFINITION 5. --- Let X be a topological space and let A be a subset of X. We call a contraction of X onto A a homotopy \(\sigma\colon X\times I\to X\) fixed on A whose origin is the identity map of X and whose term is a retraction of X onto A. If there exists a contraction of X onto A, we say that X can be contracted onto A, or again that A is a contraction of X.

Remark. --- In other words, a contraction of X onto A is a homotopy \(\sigma\colon X \times I \to X\) having the following properties:

\begin{enumerate}
\def\labelenumi{(\roman{enumi})}
\item
  \(\sigma(x,0)=x\) for all \(x\in X\);
\item
  \(\sigma(x,1)\in\mathrm{A}\) for all \(x\in\mathrm{X}\);
\item
  \(\sigma(x,t)=x\) for all \(x\in\mathbf{A}\) and every \(t\in\mathbf{I}\).
\end{enumerate}

For there to exist a contraction of X onto A, it is necessary and sufficient that the canonical injection of A into X be a homotopy equivalence of the pair (A,A) onto the pair (X,A).

Let \(a\) be a point of X. A contraction of X onto the subspace \(\{a\}\) is none other than a homotopy pointed at \(a\) connecting the map \(\mathrm{Id}_{\mathrm{X}}\) to the constant map with image \(\{a\}\). Consequently, X contracts onto \(\{a\}\) if and only if X is contractible at \(a\) (example 3 of III, p. 234).

Let X be a topological space and let A be a subset of X. Let \(\sigma\) be a contraction of X onto A. The map r from X to A defined by \(r(x) = \sigma(x,1)\) is a retraction of X onto A and \(\sigma\) is a homotopy connecting \(\mathrm{Id}_{\mathrm{X}}\) to \(r\). The relations \(r \circ i = \mathrm{Id}_{\mathrm{A}}\) and \(i \circ r = r\) then imply that the maps \(i\) and \(r\) are homotopy equivalences, inverse to each other up to homotopy.

DEFINITION 6. --- With the preceding notation, \(\sigma\) is said to be a strong contraction if, in addition, one has \(r(\sigma(x,t))=r(x)\) for all \(x\in X\) and all \(t\in I\).

Example. --- Let X be the complement of the origin in \(B_{n}\). The map \(X \times I\) to X given by \((x, t) \mapsto ((1 - t) + t \frac{1}{\|x\|})x\) is a strong contraction of X onto \(S_{n-1}\). The retraction of X onto \(S_{n-1}\) associated with it is the map given by \(x \mapsto x / \|x\|\).

Lemma 2. --- Let X be a topological space, let U be an open subset of X, and let \(\sigma\) be a strong contraction of X onto \(X-U\). Then \(\sigma(U\times I)\subset\overline{U}\). In particular, \(\sigma(U\times\{1\})\) is contained in the boundary of U.

Let \(x \in \mathrm{U}\). The set of real numbers \(t \in \mathbf{I}\) such that \(\sigma(x, t) \in \mathrm{U}\) is open in \(\mathbf{I}\) and contains 0; let \(s\) be its upper bound. We have \(\sigma(x, s) \in \overline{\mathrm{U}}\). If \(s < 1\), \(\sigma(x, s) \notin \mathrm{U}\) by definition of \(s\); the same holds if \(s = 1\), since \(\sigma(x, 1) \in \mathrm{X} - \mathrm{U}\). Consequently, \(\sigma(x, s) \in \mathrm{Fr}(\mathrm{U})\). By definition of a strong contraction, we then have \(\sigma(x, s) = \sigma(\sigma(x, s), 1) = \sigma(x, 1)\); in particular, \(\sigma(x, 1) \in \overline{\mathrm{U}}\). Consequently, \(\sigma(x, 1)\) belongs to the boundary \(\overline{\mathrm{U}} \cap (\mathrm{X} - \mathrm{U})\) of \(\mathrm{U}\).

Let \(t \in \mathbf{I}\); if \(\sigma(x,t) \notin \mathrm{U}\), we also get \(\sigma(x,t) = \sigma(\sigma(x,t),1) = \sigma(x,1)\), therefore \(\sigma(x,t) \in \overline{\mathrm{U}}\), whence the lemma.

